\subsection{Initial Conditions and Global Structure}
  \label{subsec:initial-conditions-and-global-structure}

  The framework does not postulate initial conditions in the conventional
  temporal sense.
  It assumes that~$\chi$ admits a minimal admissible ordering state,
  denoted~$\chi_0$, defining a structural boundary of admissible projected
  descriptions.
  This state does not correspond to a distinguished moment in time but to the
  earliest configurations for which an effective ordering interpretation
  becomes meaningful.

  In effective geometric regimes, the characteristic scale near~$\chi_0$
  coincides numerically with the Planck scale.
  This correspondence reflects the breakdown of projectability below this
  regime, not the presence of a fundamental cutoff or underlying discreteness.

  Cosmic history is interpreted as the progressive and irreversible growth of
  stably projected structure, accumulating away from this minimal admissible
  boundary.
  No spacetime singularity is required at the fundamental level.
  Apparent singular behavior arises only when classical notions are
  extrapolated beyond the regime in which projected configurations admit a
  stable geometric interpretation.

  \paragraph{Ontological poverty and the growth of admissible structure.}
    The minimal state~$\chi_0$ corresponds to a regime of \emph{ontological
poverty}: only a severely restricted class of simple and highly coherent
    configurations can be projected.
    As the admissibility cascade advances, the space of stably projected
    configurations expands, enabling the emergence of increasingly rich,
    localized, and hierarchical effective structures.
    The global structure of admissible descriptions is constrained by the
    ordering properties of the projective structure.

    The structural properties of~$\chi_0$ are constrained from below as well.
    A state that admitted a globally injective effective description — one in
    which every substrate configuration projected uniquely onto an observable
    — would correspond to the global minimum of the cascade and would terminate
    all temporal ordering.
    The non-commutativity $[X,\sigma(X)] = Z \neq 0$ of the admissible fibre
    generators~\cite{Beau2026m} guarantees that no such perfectly closed
    state can be admissible: every admissible projection has a non-trivial
    kernel, so the minimal state~$\chi_0$ necessarily carries an irreducible
    fibre multiplicity.
    It is this multiplicity that generates the first distinction and initiates
    the cascade without any exogenous trigger.
    The question of initial conditions thus dissolves: $\chi_0$ is not a
    fine-tuned boundary condition but the unique admissible limit of a
    projection that is structurally prevented from being perfect.
